\documentclass[10pt]{article}
\usepackage[margin=0.58in]{geometry}
\usepackage{amsmath,amssymb,array,booktabs,multicol,xcolor,enumitem}
\usepackage[T1]{fontenc}
\usepackage{lmodern}
\pagestyle{empty}
\setlength{\parindent}{0pt}
\setlength{\parskip}{2pt}
\setlength{\columnsep}{0.28in}
\setlist[itemize]{leftmargin=1.3em,itemsep=1pt,topsep=1pt}

\definecolor{deepgreen}{HTML}{244734}
\definecolor{softgreen}{HTML}{E9EFEA}
\definecolor{gold}{HTML}{9A7A38}

\newcommand{\sect}[1]{%
  \vspace{4pt}%
  \noindent\colorbox{softgreen}{\parbox{\dimexpr\linewidth-2\fboxsep\relax}{\color{deepgreen}\bfseries #1}}%
  \vspace{3pt}%
}
\newcommand{\ann}[1]{a_{\overline{#1}|}}
\newcommand{\annDue}[1]{\ddot{a}_{\overline{#1}|}}
\newcommand{\acc}[1]{s_{\overline{#1}|}}
\newcommand{\accDue}[1]{\ddot{s}_{\overline{#1}|}}
\newcommand{\PV}{\operatorname{PV}}
\newcommand{\FV}{\operatorname{FV}}

\begin{document}

\begin{center}
{\LARGE\bfseries\color{deepgreen} Financial Mathematics Formula Sheet}\\[2pt]
{\small An independent reference for interest theory and Exam FM preparation}\\[4pt]
{\footnotesize Standard notation: $i$ = effective interest rate per period, $v=(1+i)^{-1}$, $d=i/(1+i)$, and $n$ = number of periods.}
\end{center}

\begin{multicols}{2}

\sect{1. Interest Measurement}

If $a(t)$ is the accumulation function with $a(0)=1$, then an investment of $K$ at time $0$ has value $K a(t)$ at time $t$.

\[ i_{[t_1,t_2]}=\frac{a(t_2)-a(t_1)}{a(t_1)}, \qquad d_{[t_1,t_2]}=\frac{a(t_2)-a(t_1)}{a(t_2)}. \]

For one period,
\[ v=\frac{1}{1+i}=1-d, \qquad d=\frac{i}{1+i}, \qquad i=\frac{d}{1-d}. \]

\textbf{Simple interest:}
\[ a(t)=1+it. \]

\textbf{Simple discount:}
\[ a(t)=\frac{1}{1-dt}. \]

\textbf{Compound interest:}
\[ a(t)=(1+i)^t, \qquad v^t=(1+i)^{-t}. \]

For nominal rates convertible $m$ times per year,
\[ 1+i=\left(1+\frac{i^{(m)}}{m}\right)^m, \qquad 1-d=\left(1-\frac{d^{(m)}}{m}\right)^m. \]

For a constant force of interest $\delta$,
\[ 1+i=e^{\delta}, \qquad \delta=\ln(1+i). \]

For a time-varying force $\delta_t$,
\[ \delta_t=\frac{a'(t)}{a(t)}, \qquad a(t)=\exp\!\left(\int_0^t \delta_s\,ds\right). \]

\sect{2. Equations of Value}

Move every cash flow to one common comparison date. If cash flow $C_k$ occurs at time $t_k$, then at time $0$,
\[ \PV=\sum_k C_k v^{t_k}. \]

At time $T$,
\[ \FV_T=\sum_k C_k(1+i)^{T-t_k}. \]

The yield rate or internal rate of return is a rate $i$ satisfying
\[ \sum_k C_k(1+i)^{-t_k}=0, \]
where inflows and outflows are assigned opposite signs.

\sect{3. Level Annuities}

For payments of $1$ at times $1,2,\ldots,n$,
\[ \ann{n}=\frac{1-v^n}{i}, \qquad \acc{n}=\frac{(1+i)^n-1}{i}. \]

For payments of $1$ at times $0,1,\ldots,n-1$,
\[ \annDue{n}=(1+i)\ann{n}=\frac{1-v^n}{d}, \]
\[ \accDue{n}=(1+i)\acc{n}. \]

Perpetuities:
\[ \ann{\infty}=\frac{1}{i}, \qquad \annDue{\infty}=\frac{1}{d}. \]

A deferred annuity-immediate with first payment after $k$ deferral periods has present value
\[ {}_{k|}\ann{n}=v^k\ann{n}. \]

For level payment $R$, multiply the appropriate annuity factor by $R$.

\sect{4. Varying Annuities}

For payments $1,2,\ldots,n$ at times $1,2,\ldots,n$,
\[ (I\!a)_{\overline{n}|}=\frac{\annDue{n}-nv^n}{i}. \]

For payments $n,n-1,\ldots,1$,
\[ (D\!a)_{\overline{n}|}=\frac{n-\ann{n}}{i}. \]

For payments $1,(1+g),\ldots,(1+g)^{n-1}$,
\[ \PV=\frac{1-\left(\frac{1+g}{1+i}\right)^n}{i-g}, \qquad i\neq g. \]

For an infinite geometric payment stream beginning with $1$ at time $1$,
\[ \PV=\frac{1}{i-g}, \qquad i>g. \]

\sect{5. Payment Frequency}

If the annual effective rate is $i$, define
\[ i^{(m)}=m\big((1+i)^{1/m}-1\big), \]
\[ d^{(m)}=m\big(1-(1+i)^{-1/m}\big). \]

For $m$ payments per year of amount $1/m$ over $n$ years,
\[ \ann{n}^{(m)}=\frac{1-v^n}{i^{(m)}}=\ann{n}\frac{i}{i^{(m)}}, \]
\[ \annDue{n}^{(m)}=\frac{1-v^n}{d^{(m)}}=\ann{n}\frac{i}{d^{(m)}}. \]

For a continuous payment rate $c(t)$,
\[ \PV=\int_0^n c(t)v(t)\,dt. \]
If the force is constant, $v(t)=e^{-\delta t}$.

\sect{6. Loans}

For a loan of amount $L$ repaid by $n$ level payments $R$,
\[ L=R\ann{n}, \qquad R=\frac{L}{\ann{n}}. \]

Immediately after payment $t$, the outstanding balance may be found prospectively:
\[ B_t=R\ann{n-t}, \]
or retrospectively:
\[ B_t=L(1+i)^t-R\acc{t}. \]

For nonlevel payments,
\[ B_{t}=B_{t-1}(1+i)-R_t. \]

Interest and principal in payment $t$ are
\[ I_t=iB_{t-1}, \qquad P_t=R_t-I_t. \]

Principal repaid from time $r$ through time $s$ is the decrease in balance over that interval.

\sect{7. Bonds}

Let $F$ be face value, $r$ the coupon rate per coupon period, $C$ the redemption value, $j$ the yield rate per coupon period, and $n$ the number of remaining coupons. Coupon amount is $Fr$.

Bond price:
\[ P=Fr\,\ann{n}\big|_{i=j}+Cv_j^n. \]

Equivalent premium/discount form:
\[ P-C=(Fr-Cj)\ann{n}\big|_{i=j}. \]

Thus $P>C$ when $Fr>Cj$ (premium) and $P<C$ when $Fr<Cj$ (discount).

Book value immediately after coupon $t$:
\[ B_t=Fr\,\ann{n-t}\big|_{i=j}+Cv_j^{n-t}. \]

Recursive book value:
\[ B_t=B_{t-1}(1+j)-Fr. \]

Interest earned during period $t$ is $jB_{t-1}$; the difference between the coupon and interest is the premium amortization or discount accumulation.

\sect{8. Spot and Forward Rates}

If $s_n$ is the $n$-period spot rate,
\[ v(0,n)=(1+s_n)^{-n}. \]

If $f_{n,n+m}$ is the $m$-period forward rate beginning at time $n$,
\[ (1+s_{n+m})^{n+m}=(1+s_n)^n(1+f_{n,n+m})^m. \]

For one-period forward rates,
\[ (1+s_n)^n=\prod_{k=1}^{n}(1+f_{k-1,k}). \]

Value any deterministic cash flow stream using spot rates:
\[ P=\sum_{t} \frac{C_t}{(1+s_t)^t}. \]

\sect{9. Duration and Convexity}

For price
\[ P(i)=\sum_t C_t(1+i)^{-t}, \]
Macaulay duration is
\[ D_M=\frac{\sum_t t\,C_t(1+i)^{-t}}{P}. \]

Modified duration is
\[ D_{\mathrm{mod}}=\frac{D_M}{1+i}=-\frac{1}{P}\frac{dP}{di}. \]

A useful convexity measure is
\[ \mathcal C=\frac{1}{P}\frac{d^2P}{di^2}=\frac{\sum_t t(t+1)C_t(1+i)^{-t-2}}{P}. \]

For a small change $\Delta i$,
\[ \frac{\Delta P}{P}\approx -D_{\mathrm{mod}}\Delta i+\frac12\mathcal C(\Delta i)^2. \]

\sect{10. Matching and Immunization}

\textbf{Exact matching:} choose asset cash flows so that the asset cash flow at each relevant date equals the liability cash flow at that date.

\textbf{First-order immunization at rate $i$:}
\begin{itemize}
\item present value of assets = present value of liabilities;
\item modified duration of assets = modified duration of liabilities.
\end{itemize}

A standard Redington condition additionally requires asset convexity to exceed liability convexity at the matching rate.

\sect{11. Quick Identities}

\[ \annDue{n}-\ann{n}=1-v^n, \qquad \accDue{n}-\acc{n}=(1+i)^n-1. \]

\[ \ann{n}=v\annDue{n}, \qquad \acc{n}=(1+i)^n\ann{n}. \]

For a level annuity, always check whether the valuation date is \emph{before the first payment} (immediate) or \emph{at the first payment} (due).

\end{multicols}

\end{document}
